%% parts/p06/infinitary-combinatorics.tex --- Phase-4 P6 depth fill
%%
%% Register coverage: D-P7-005, D-P7-009, D-P7-010;
%% N-a-037--N-a-041; apparatus A-064.  The iterated-forcing consistency
%% proof for Martin's axiom remains excluded by N-x-003.

\chapter{Infinitary Combinatorics}
\label{sec:p06-infinitary-combinatorics}
\chaptermark{Infinitary Combinatorics}

This chapter fixes the hypotheses behind the combinatorial principles used
later in set theory.  Regularity and uncountability are stated rather than
left implicit, and principles with similar-looking symbols are kept distinct
in standing.  Standard background and the calibrated implications below are
from \cite{Jech2003,Kunen2011}.

\section{Clubs, stationary sets, and pressing down}
\label{sec:p06-club-stationary}

\begin{hbconvention}{def:p06-regular-uncountable-context}{Regular-uncountable context}
Throughout this section, \(\kappa\) is a regular uncountable cardinal unless
a statement explicitly says otherwise.  Here regularity means that every set
of fewer than \(\kappa\) ordinals below \(\kappa\) has supremum below
\(\kappa\).  This local exact clause avoids a forward definition; the general
cofinality entry is Definition~\ref{def:p06-cofinality-regular-singular}.
\end{hbconvention}

\begin{hbdefinition}{def:p06-club-stationary}{Club and stationary subsets}
Let \(C,S\subseteq\kappa\).
\begin{enumerate}
  \item \(C\) is \emph{unbounded} in \(\kappa\) iff
  \[
    \forall\alpha<\kappa\,\exists\beta\in C\;(\alpha<\beta).
  \]
  \item \(C\) is \emph{closed} in \(\kappa\) iff, for every nonzero limit
        ordinal \(\delta<\kappa\), equality
        \(\sup(C\cap\delta)=\delta\) implies \(\delta\in C\).
  \item \(C\) is a \emph{club} subset of \(\kappa\) iff it is closed and
        unbounded.
  \item \(S\) is \emph{stationary} in \(\kappa\) iff
        \(S\cap C\ne\varnothing\) for every club \(C\subseteq\kappa\).
\end{enumerate}
The hypotheses on \(\kappa\) belong to the definition's active scope; the
same words over singular cardinals require separate conventions.
\index{closed unbounded set}\index{club set}\index{stationary set}
\end{hbdefinition}

\begin{hbdefinition}{def:p06-diagonal-intersection}{Diagonal intersection}
For a sequence \(\langle C_\xi:\xi<\kappa\rangle\) of subsets of \(\kappa\),
its \emph{diagonal intersection} is
\[
  \mathop{\triangle}_{\xi<\kappa}C_\xi
  \defeq
  \{\alpha<\kappa:\forall\xi<\alpha\;(\alpha\in C_\xi)\}.
\]
The \emph{club filter} on \(\kappa\) consists of the subsets of \(\kappa\)
that contain a club.  Its dual ideal consists of the nonstationary subsets.
A filter on \(\kappa\) is \emph{\(\kappa\)-complete} iff it is closed under
intersections of nonempty families of fewer than \(\kappa\) members, and it is
\emph{normal} iff it is closed under the displayed diagonal intersections of
\(\kappa\)-sequences of its members.
Here \(\mathop{\triangle}\) is a locally declared operation on a
\(\kappa\)-sequence of sets, not a formula connective.
\index{diagonal intersection}\index{club filter}
\index{nonstationary ideal}
\end{hbdefinition}

\begin{hbtheorem}{thm:p06-club-normality}{cited=Jech2003}{Closure and normality of the club filter}
Let \(0<\lambda<\kappa\).  The intersection of any
\(\lambda\)-sequence of club subsets of \(\kappa\) is club.  If
\(\langle C_\xi:\xi<\kappa\rangle\) is a sequence of club subsets, then
\[
  \mathop{\triangle}_{\xi<\kappa}C_\xi
\]
is club.  Consequently the club filter is \(\kappa\)-complete and normal.
\end{hbtheorem}

\begin{hbtheorem}{thm:p06-fodor}{cited=Fodor1956}{Fodor's pressing-down lemma}
Let \(S\subseteq\kappa\) be stationary and let \(f:S\to\kappa\) be
\emph{regressive}: \(f(\alpha)<\alpha\) for every nonzero
\(\alpha\in S\).  Then for some \(\beta<\kappa\), the fiber
\[
  \{\alpha\in S:f(\alpha)=\beta\}
\]
is stationary.  The conclusion is stationary, not merely uncountable
\cite{Fodor1956}.
\index{Fodor's lemma}\index{pressing-down lemma}
\end{hbtheorem}

\section{Delta systems and partition relations}
\label{sec:p06-delta-ramsey}

\begin{hbdefinition}{def:p06-delta-system}{Delta system}
A family \(\mathcal B\) of sets is a \emph{delta system} with
\emph{root} \(R_{\mathcal B}\) iff
\[
  x\cap y=R_{\mathcal B}
  \qquad
  \text{whenever }x,y\in\mathcal B\text{ and }x\ne y.
\]
The root is part of the conclusion; it is not assumed in advance.
\index{Delta-system lemma}\index{root!Delta system}
\end{hbdefinition}

\begin{hbtheorem}{thm:p06-delta-system}{proved=annex:p6.comb.1}{Delta-system lemma}
In \(\mathsf{ZFC}\), every uncountable family of finite sets has an
uncountable subfamily that is a delta system with a finite root.  No
generalization to families of sets of size below an arbitrary \(\kappa\) is
asserted here; such versions require additional hypotheses
\cite{Shanin1946,Jech2003}.
\end{hbtheorem}

\begin{hblemma}{thm:p06-add-ccc}{proved=annex:p6.comb.2}{The Cohen poset is ccc}
For every infinite cardinal \(\kappa\),
\(\operatorname{Add}(\omega,\kappa)\) has the countable chain condition.
\end{hblemma}

Its complete proof is
Annex~P6.COMB.2.

\begin{hbdefinition}{def:p06-partition-arrow}{Partition-arrow notation}
For cardinals \(\kappa,\lambda,\mu\) and a positive integer \(\ell\), write
\[
  \kappa\longrightarrow(\lambda)^\ell_\mu
\]
iff every coloring \(c:[\kappa]^\ell\to\mu\) has a set
\(H\subseteq\kappa\) with \(\abs H=\lambda\) such that \(c\) is constant on
\([H]^\ell\).  Here \([X]^\ell\) is the set of \(\ell\)-element subsets of \(X\),
and the displayed arrow denotes a partition relation, not implication in an
object language.
\index{partition relation}\index{arrow notation!partition calculus}
\end{hbdefinition}

\begin{hbtheorem}{thm:p06-infinite-ramsey}{cited=Ramsey1930}{Infinite Ramsey theorem}
For every positive finite \(\ell\) and every positive finite \(k\),
\[
  \omega\longrightarrow(\omega)^\ell_k.
\]
Equivalently, every finite coloring of the \(\ell\)-element subsets of
\(\omega\) has an infinite homogeneous set \cite{Ramsey1930}.
\end{hbtheorem}

\begin{hbtheorem}{thm:p06-erdos-rado}{cited=ErdosRado1956}{Erdos--Rado theorem}
For an infinite cardinal \(\kappa\), define the finite beth iteration by
\[
  \beth_0(\kappa)=\kappa,
  \qquad
  \beth_{n+1}(\kappa)=2^{\beth_n(\kappa)}.
\]
Then for every \(n<\omega\),
\[
  \beth_n(\kappa)^+
  \longrightarrow
  (\kappa^+)^{n+1}_\kappa.
\]
This is the ordinary, unpolarized partition relation; no singular or
polarized strengthening is included \cite{ErdosRado1956}.
\index{Erdos--Rado theorem}\index{beth iteration}
\end{hbtheorem}

\begingroup\small
\begin{longtable}{@{}p{0.24\textwidth}p{0.38\textwidth}p{0.30\textwidth}@{}}
\caption{Partition-relation register for the infinitary-combinatorics
entry.}\label{tab:p06-partition-relations}\\
\toprule
Item & Exact form & Scope and status \\
\midrule
\endfirsthead
\toprule
Item & Exact form & Scope and status \\
\midrule
\endhead
Arrow notation &
\(\kappa\longrightarrow(\lambda)^\ell_\mu\) &
Colorings \([\kappa]^\ell\to\mu\); a homogeneous set has cardinality
exactly \(\lambda\). \\
\addlinespace
Infinite Ramsey &
\(\omega\longrightarrow(\omega)^\ell_k\) &
Finite positive \(\ell,k\); cited as Theorem~\ref{thm:p06-infinite-ramsey}. \\
\addlinespace
Erdos--Rado &
\(\beth_n(\kappa)^+\longrightarrow(\kappa^+)^{n+1}_\kappa\) &
Infinite \(\kappa\), finite \(n\); cited as
Theorem~\ref{thm:p06-erdos-rado}. \\
\bottomrule
\end{longtable}
\endgroup

\section{Trees and the Suslin hypothesis}
\label{sec:p06-trees-suslin}

\begin{hbdefinition}{def:p06-omega-one-trees}{Trees, levels, and branches}
A \emph{tree} is a partially ordered set \((\mathcal T,<_{\mathcal T})\) such
that the strict predecessors of every node are well-ordered by
\(<_{\mathcal T}\).  The height of a
node is the order type of its predecessor set, and
\[
  \mathcal T_\alpha\defeq
  \{\mathbf t\in\mathcal T:
    \operatorname{ht}_{\mathcal T}(\mathbf t)=\alpha\}
\]
is the \(\alpha\)-th level.  A \emph{branch} is a maximal chain.  If
\(\lambda\) is a regular uncountable cardinal, a \emph{\(\lambda\)-tree}
has height \(\lambda\) and satisfies \(\abs{\mathcal T_\alpha}<\lambda\) for every
\(\alpha<\lambda\).  A \emph{\(\lambda\)-Aronszajn tree} is a
\(\lambda\)-tree with no branch of order type \(\lambda\).  Thus an
\(\omega_1\)-tree has countable levels, and an unqualified
\emph{Aronszajn tree} means an \(\omega_1\)-Aronszajn tree.
\index{tree!set-theoretic}\index{level!of a tree}
\index{Aronszajn tree}
\end{hbdefinition}

The choice-free binary-tree result needed by this chapter is already
Theorem~\ref{thm:p06-konig-infinity}; its restriction to
\(\mathcal T\subseteq2^{<\omega}\) is deliberate.  No assertion about arbitrary
finitely branching trees is substituted for it.

\begin{hbtheorem}{thm:p06-aronszajn-existence}{cited=Kurepa1935}{Existence of an Aronszajn tree}
\(\mathsf{ZFC}\) proves that an Aronszajn tree exists.  The cited construction
verifies separately that every level is countable and that there is no branch
of order type \(\omega_1\); it is not replaced here by a one-line construction
pointer \cite{Kurepa1935,Todorcevic1984Trees}.
\end{hbtheorem}

\begin{hbmethod}{def:p06-aronszajn-construction}{A special Aronszajn-tree construction}
Build levels \(U_\alpha\) for \(\alpha<\omega_1\) from strictly increasing
sequences of rationals in \((0,1)\), ordered by end-extension, while
maintaining that every node can be extended to any higher countable level
with any prescribed rational supremum above its current supremum; put
\(\sup\varnothing=0\) for the empty node.  At a
successor level, extend each node by each admissible rational.  At a countable
limit \(\alpha\), the earlier tree is countable.  For every earlier node
\(s\) and rational
\(\sup\operatorname{ran}(s)<q<1\), choose a cofinal
\(\omega\)-sequence in \(\alpha\), extend \(s\) along it with rational values
increasing to \(q\), and put the union sequence on level \(\alpha\).
There are only countably many pairs \((s,q)\), so every level is countable
and the extension invariant continues.

For a node \(s\), the rational \(\sup\operatorname{ran}(s)\) strictly
increases under proper end-extension.  A branch of order type \(\omega_1\)
would therefore inject \(\omega_1\) into \(\mathbb Q\), which is impossible.
The construction has a nonempty countable level at every countable ordinal
and no cofinal branch, so it yields the cited Aronszajn tree.  The method
records the level recursion and branch obstruction; full bookkeeping remains
with \cite{Kurepa1935,Todorcevic1984Trees}.
\end{hbmethod}

\begin{hbdefinition}{def:p06-suslin-hypothesis}{Suslin trees and the Suslin hypothesis}
A \emph{Suslin tree} is an \(\omega_1\)-tree with neither an uncountable
chain nor an uncountable antichain.  The \emph{Suslin hypothesis}
\((\mathsf{SH})\) asserts that no Suslin tree exists.  A \emph{Suslin line}
is a Dedekind-complete dense linear order without endpoints whose order
topology is ccc but not separable.  Under the standard tree--line
correspondence, \(\mathsf{SH}\) is equivalent to the nonexistence of such a
line \cite{Todorcevic1984Trees}.
\index{Suslin tree}\index{Suslin hypothesis}\index{Suslin line}
\end{hbdefinition}

\begin{hbtheorem}{thm:p06-sh-relative-status}{cited=Jech2003}{Relative status of the Suslin hypothesis}
Relative to the consistency of \(\mathsf{ZFC}\), both
\(\mathsf{ZFC}+\mathsf{SH}\) and
\(\mathsf{ZFC}+\neg\mathsf{SH}\) are consistent.  The negative direction is
obtained from the diamond principle in the constructible universe; the
positive direction follows from a forcing model of
\(\mathsf{MA}_{\aleph_1}\).  These relative-consistency constructions are
cited, not reproduced locally.
\end{hbtheorem}

\section{Diamond, square, and Martin's axiom}
\label{sec:p06-combinatorial-principles}

\begin{hbdefinition}{def:p06-diamond}{The diamond principle at omega-one}
The principle \(\diamondsuit_{\omega_1}\) asserts that there is a sequence
\(\langle A_\alpha:\alpha<\omega_1\rangle\) such that
\(A_\alpha\subseteq\alpha\) and, for every \(A\subseteq\omega_1\),
\[
  \{\alpha<\omega_1:A\cap\alpha=A_\alpha\}
\]
is stationary in \(\omega_1\).  The suit-shaped \(\diamondsuit\) denotes a
set-theoretic guessing principle, not the modal possibility operator.
\index{diamond principle@\(\diamondsuit\) principle}
\end{hbdefinition}

\begin{hbtheorem}{thm:p06-diamond-suslin}{cited=Jensen1972FineStructure}{Diamond produces a Suslin tree}
The principle \(\diamondsuit_{\omega_1}\) implies that an
\(\omega_1\)-Suslin tree exists, and hence implies
\(\neg\mathsf{SH}\).  In particular, the constructible universe satisfies
the relevant diamond principle \cite{Jensen1972FineStructure}.
\end{hbtheorem}

\begin{hbdefinition}{def:p06-square}{The square principle}
For a nonzero limit ordinal \(\alpha\), a set \(C\subseteq\alpha\) is
\emph{club in \(\alpha\)} when it is unbounded in \(\alpha\) and, for every
nonzero limit \(\beta<\alpha\), unboundedness of \(C\cap\beta\) in \(\beta\)
implies \(\beta\in C\).  This local ordinal convention is needed because the
indices below need not be regular cardinals.

For an infinite cardinal \(\kappa\), the principle \(\square_\kappa\)
asserts that there is a sequence
\[
  \langle C_\alpha:
  \alpha<\kappa^+\text{ and }\alpha\text{ is a nonzero limit ordinal}\rangle
\]
such that:
\begin{enumerate}
  \item \(C_\alpha\) is club in \(\alpha\) and
        \(\operatorname{otp}(C_\alpha)\leq\kappa\);
  \item if \(\beta\) is a limit point of \(C_\alpha\), then
        \(C_\beta=C_\alpha\cap\beta\); and
  \item there is no club \(C\subseteq\kappa^+\) such that
        \(C\cap\alpha=C_\alpha\) at every nonzero limit point
        \(\alpha\) of \(C\).
\end{enumerate}
The last clause is the \emph{no-thread condition}; coherent clubs without it
do not state this square principle.  The square-shaped glyph here names a
set-theoretic sequence principle, not the modal necessity operator.
\index{square principle@\(\square\) principle}\index{thread!square sequence}
\end{hbdefinition}

\begin{hbtheorem}{thm:p06-square-aronszajn}{cited=Jensen1972FineStructure}{Square produces an Aronszajn tree}
For every infinite cardinal \(\kappa\), the full principle
\(\square_\kappa\) implies that a \(\kappa^+\)-Aronszajn tree exists.  This
is an application of the no-thread clause, not merely of coherence
\cite{Jensen1972FineStructure}.
\end{hbtheorem}

\begin{hbdefinition}{def:p06-martins-axiom}{Martin's axiom}
Let \(\mathfrak c\defeq2^{\aleph_0}\).  For an infinite cardinal \(\kappa\),
\(\mathsf{MA}_\kappa\) asserts: whenever \(\mathbb P\) is ccc and
\(\mathcal D\) is a family of at most \(\kappa\) dense subsets of
\(\mathbb P\), some poset filter \(G\subseteq\mathbb P\) meets every member
of \(\mathcal D\).  The terms ccc, dense, and poset filter are exactly those
of Definitions~\ref{def:p06-ccc} and \ref{def:p06-poset-filter}.
Martin's axiom \(\mathsf{MA}\) is the scheme
\[
  \mathsf{MA}_\kappa\qquad(\kappa<\mathfrak c).
\]
It is a forcing axiom about filters meeting specified dense sets, not an
assertion that a generic filter over the ambient universe exists.
\index{Martin's axiom}
\end{hbdefinition}

\begin{hbtheorem}{thm:p06-ma-suslin}{cited=MartinSolovay1970}{Martin's axiom and Suslin trees}
The axiom \(\mathsf{MA}_{\aleph_1}\) implies the Suslin hypothesis.  Hence
\(\mathsf{MA}+\neg\mathsf{CH}\) implies \(\mathsf{SH}\), since then
\(\aleph_1<\mathfrak c\) \cite{MartinSolovay1970,Jech2003}.
\end{hbtheorem}

\begin{proofboundary}{def:p06-ma-proof-boundary}{Martin's-axiom proof boundary}
The consistency proof for \(\mathsf{MA}+\neg\mathsf{CH}\) requires an
iteration with bookkeeping and preservation of the ccc.  That specialist
forcing technology is outside this chapter's scope.  The chapter states
the axiom and its calibrated consequences but supplies no annex item for the
iteration.
\end{proofboundary}

\begingroup\footnotesize
\begin{longtable}{@{}p{0.19\textwidth}p{0.48\textwidth}p{0.25\textwidth}@{}}
\caption{Combinatorial principles and their implication structure.}\label{tab:p06-combinatorial-principles}\\
\toprule
Principle & Exact hypothesis and conclusion & Consequence and evidence \\
\midrule
\endfirsthead
\toprule
Principle & Exact hypothesis and conclusion & Consequence and evidence \\
\midrule
\endhead
Club / stationary &
\(C\subseteq\kappa\) is closed and unbounded; \(S\subseteq\kappa\) meets
every club.  Here \(\kappa\) is regular uncountable. &
Definitions~\ref{def:p06-club-stationary} and
\ref{def:p06-diagonal-intersection}; definitional rows. \\
\addlinespace
Diagonal club &
The diagonal intersection of a \(\kappa\)-sequence of clubs is club. &
Normality of the club filter; cited to \textcite{Jech2003} in
Theorem~\ref{thm:p06-club-normality}. \\
\addlinespace
Fodor &
A regressive map on a stationary subset of regular uncountable \(\kappa\)
has a stationary constant fiber. &
Cited to \textcite{Fodor1956} in Theorem~\ref{thm:p06-fodor}; regularity and
uncountability are essential. \\
\addlinespace
Delta system &
Every uncountable family of finite sets has an uncountable fixed-root
subfamily. &
Proved in Annex~P6.COMB.1; the separately proved Cohen application gives
\(\operatorname{Add}(\omega,\kappa)\) ccc in Annex~P6.COMB.2. \\
\addlinespace
Ramsey / Erdos--Rado &
\(\omega\longrightarrow(\omega)^\ell_k\) for positive finite \(\ell,k\); and
\(\beth_n(\kappa)^+\longrightarrow(\kappa^+)^{n+1}_\kappa\) for infinite
\(\kappa\) and finite \(n\). &
Cited to \textcite{Ramsey1930} and \textcite{ErdosRado1956}. \\
\addlinespace
\(\diamondsuit_{\omega_1}\) &
A sequence guesses every \(A\subseteq\omega_1\) correctly on a stationary
set of initial segments. &
It implies \(\neg\mathsf{SH}\) via a Suslin tree; cited to
\textcite{Jensen1972FineStructure}. \\
\addlinespace
\(\square_\kappa\), \(\kappa\) infinite &
Coherent clubs of order type at most \(\kappa\), with no club thread. &
It implies the existence of a \(\kappa^+\)-Aronszajn tree; cited to
\textcite{Jensen1972FineStructure}. \\
\addlinespace
\(\mathsf{MA}_\kappa\), \(\kappa\) infinite &
Every family of at most \(\kappa\) dense sets in a ccc poset is met by one
poset filter. &
The \(\aleph_1\) instance implies \(\mathsf{SH}\), cited to
\textcite{MartinSolovay1970}; the iteration proof is outside this chapter. \\
\bottomrule
\end{longtable}
\endgroup

\paragraph{Source disposition.}
The historical v1.1 Part~VII source contains no chapter on clubs and stationary sets,
Delta systems, partition relations, Aronszajn or Suslin trees, diamond,
square, or Martin's axiom; these are separately sourced editorial additions.
Accordingly no inherited proof has been compressed in this chapter.  Two existing
results are deliberately not duplicated: binary K\"onig is
Theorem~\ref{thm:p06-konig-infinity} with Annex~P6.CHOICE.2, and the Cohen-ccc
application is Theorem~\ref{thm:p06-add-ccc} with Annex~P6.COMB.2.  The
iterated-forcing consistency proof for Martin's axiom remains outside scope.
